\documentclass[11pt,a4paper]{article}
\usepackage[margin=26mm]{geometry}
\usepackage{amsmath,amssymb,amsthm,microtype}
\usepackage[hidelinks]{hyperref}
\newcommand{\E}{\mathbb E}\newcommand{\Pp}{\mathbb P}
\newtheorem{theorem}{Theorem}\newtheorem{lemma}{Lemma}
\title{Direct count comparison with the physical variance allocation}
\author{Research note, 28 October 2026}\date{}
\begin{document}\maketitle

The difficult negative phases can be handled without reducing the conditional
support below $3/5$. The useful information is asymmetric: a negative noise
variable can make a large negative jump, whereas a positive noise variable
has negative jumps of size at most $9/125$. We retain this distinction for the
first incoming count, and retain the same conditional third moments in the
outgoing count and in the cubic penalty. A small closed certificate completes
the resulting elementary comparison. The certificate has 291 cells; it does
not discretize the sample size or the marginal distribution.

\begin{theorem}\label{thm:comparison}
Let $n\ge42$, $2/5\le p\le53/100$, $q=1-p$, and $Y=B+U$, where
$B$ is Bernoulli with parameter $p$ and $\E(U\mid B)=0$. Write $U_j$
for the conditional noise given $B=j$, and assume
\[
 U_0\in[-A,a_0],\qquad U_1\in[-a_1,A],\qquad
 A=\frac35,\quad a_0,a_1\ge0,\quad a_0+a_1\le d=\frac9{125}.
\]
Set $v=pq$, $s=\E U^2$, $\lambda=ns$, $\ell=(v+s)^{-1/2}$,
$X=\ell(Y-p)$, $\rho=\E|X|^3$ and $\rho_B=(1-2v)/\sqrt v$.
Let $S_n$ be the sum of $n$ independent copies of $X$. Suppose
\[
 \frac3{100}\le\lambda\le\frac3{40},\qquad |H|\le\frac{46}{25},
 \qquad np+H/\ell=k-r,\quad k\in\mathbb Z,\quad0\le r\le\frac3{20}.
\]
For every $C\ge2/5$, with $\delta=k-np$ and $F_B$ the distribution
function of $\operatorname{Bin}(n,p)$,
\begin{align*}
\Delta_C={}&\sqrt n\{\Pp(S_n\le H)-\Phi(H/\sqrt n)\}-C\rho\\
 &-\sqrt n\{F_B(k)-\Phi(\delta/\sqrt{nv})\}+C\rho_B<0.
\end{align*}
\end{theorem}

This statement concerns actual iid laws, including atom collisions and zero
variance in either conditional law. It does not assume a least counterexample,
contact, harmonic reserves, half-amplitude confinement, or marked crossings.
Other phases and smaller noise require their separate comparison theorems.
The present assertion alone is not a proof of an unrestricted Berry--Esseen
bound, nor does it establish the acquisition of its strip hypotheses.

\section{The count identity and its normalization}
Write $v_j=\E U_j^2$, $\eta=1/n$, and introduce the \emph{physical} allocation
\[
 w=\frac{pv_1}{s},\qquad \Lambda_0=nqv_0=\lambda(1-w),\qquad
 \Lambda_1=npv_1=\lambda w.
\]
Let $W_j$ be a sum of $j$ independent variables with law $U_1$ and
$n-j$ with law $U_0$, and put $\pi_j=\binom njp^jq^{n-j}$.
With $u=-r$, conditioning gives exactly
\begin{align*}
 \Pp(S_n\le H)-F_B(k)
 ={}&-\sum_{i=0}^{k}\pi_{k-i}\Pp(W_{k-i}>u+i)\\
 &+\sum_{i=1}^{n-k}\pi_{k+i}\Pp(W_{k+i}\le u-i).
\end{align*}
The outgoing event is strict and the incoming event closed. Put
\[
 m=\sqrt n\pi_k,\qquad
 f(u)=\frac{\Phi((\delta+u)/\sqrt{n(v+s)})-
                   \Phi(\delta/\sqrt{n(v+s)})}{\pi_k}.
\]
If $\pi_{k+i}/\pi_k\le R$ and the sum of incoming probabilities is at
most $\mathcal T$, it is sufficient to prove positivity of
\begin{equation}\label{eq:margin}
 m\{\Pp(W_k>-r)+f(-r)-R\mathcal T\}
       +C(\rho-\rho_B)-\frac74(\Xi+r)\lambda\eta,
 \qquad \Xi=\frac{231}{250}.
\end{equation}
Indeed the final term bounds the unfavorable normal variance correction:
differentiating the normal distribution function gives at most
$(1/5)v^{-3/2}\delta_+s<(7/4)(\Xi+r)s$.

Here and below $v_*:=6/25\le v$, $\eta\le1/42$ and
\[
 \sqrt{v+s}<\frac{251}{500},\quad
 \xi:=H\sqrt{v+s},\quad |\xi|<\Xi,\quad\delta=\xi+r.
\]
In particular $1\le k\le n-1$ and $9/25<k/n<23/40$.
The finite Stirling estimate, with Taylor expansion around $p$, gives
\begin{equation}\label{eq:stirling}
 m=\frac{\phi(0)}{\sqrt v}
 \exp\{-\delta^2/(2nv)+e\},\qquad
 n|e|\le Z(D):=\frac{5D}{8}+\frac38+\frac{D^3}{42}
 \quad (|\delta|\le D\le\Xi+3/20).
\end{equation}
For clarity, the constants are respectively the prefactor error, factorial
remainder, and entropy remainder. On the segment joining $p$ to $k/n$,
$|D'''(z\Vert p)|\le6$, $|(\log(z(1-z)))'|\le5/4$, and
$1/(12z(1-z))<3/8$. Applying
$0<\log(j!)-(j+1/2)\log j+j-\tfrac12\log(2\pi)<1/(12j)$
gives the errors $D^3/n^2$, $5D/(8n)$ and $3/(8n)$.
A proof of this finite factorial formula is retained in
\texttt{manuscript/arxiv-v1/appendices/count-comparison.tex}
(the paragraph beginning ``For completeness, the finite Stirling formula
used here is''). No asymptotic remainder in the noise is used.

On a closed cell $[l,L]\times[r_-,r_+]\times[w_-,w_+]$, let
$D=\Xi+r_+$. Using $19943/50000<\phi(0)<2/5$ and
$e^{-x}\ge1-x$ gives
\begin{align}
 m_-&=2\frac{19943}{50000}
       \left[1-\frac{D^2/(2v_*)+Z(D)}{42}\right]\le m\le\frac{17}{20}=:m_+,
       \label{eq:mass}\\
 R&=1+\frac{\max(0,53/100+D-1)}
              {42v_*-(3/5)\max(0,D-1)}.\label{eq:ratio}
\end{align}
For the upper mass bound apply \eqref{eq:stirling} at a mode, whose
displacement from $np$ is at most one, and use
$e^{Z(1)/42}\le(1-Z(1)/42)^{-1}$. For \eqref{eq:ratio}, the first
exact ratio is $1+(p-\delta-1)/(nv+q(\delta+1))$; all subsequent
ratios are at most one because $53/100+D<2$.
All denominators and the bounds in this paragraph are checked rationally.

For fixed count and fixed $\delta$, the derivative of $\log f'(t)$ with
respect to $s$ is
\[
 \frac{(\delta+t)^2-n(v+s)}{2n(v+s)^2}<0,
 \qquad -r\le t\le0.
\]
Indeed $|\delta+t|\le\Xi+r$ and $42v_*>(\Xi+3/20)^2$.
At $s=0$, \eqref{eq:stirling} therefore gives
\[
 \log f'(t)\le\eta P(r),\qquad
 P(r)=\frac{2\Xi r+r^2}{2v_*}+Z(\Xi+r).
\]
Let $E=42(\exp^+(P(r_+)/42)-1)$, where $\exp^+$ is the rational
upper enclosure described below. Convexity of the exponential on
$0\le\eta\le1/42$ implies
\begin{equation}\label{eq:phase}
                       f(-r)\ge-r(1+E\eta).
\end{equation}
This keeps the full $1/n$ scaling instead of introducing a constant error.

\section{The first incoming count remembers the allocation}
We use the quadratic-hinge comparison of Pinelis,
\emph{On the Bennett--Hoeffding inequality},
\href{https://arxiv.org/abs/0902.4058}{arXiv:0902.4058},
Proposition 2.8 and Corollary 2.11. Independent centered variables bounded
above by $b>0$, with total variance at most $V$, have sum dominated in every
quadratic hinge $(x-y)_+^2$ by $b(P_{V/b^2}-V/b^2)$.
It applies with an upper bound on variance, not only its exact value.
Conditioning on independent added variables preserves the comparison.
Consequently the two groups in $-W_j$ may be replaced successively by
independent centered Poisson variables with jump sizes $A$ and $d$.

\begin{lemma}\label{lem:monotone}
For fixed $y$, the function
\[
 G(T,w;y)=\E\left[
 A\left(P_{T(1-w)/A^2}-\frac{T(1-w)}{A^2}\right)
 +d\left(P_{Tw/d^2}-\frac{Tw}{d^2}\right)-y\right]_+^2
\]
is nondecreasing in $T\ge0$ and nonincreasing in $0\le w\le1$.
\end{lemma}
\begin{proof}
Increasing $T$ adds independent centered Poisson increments, so conditional
Jensen applies. To compare jump sizes at fixed variance $V$, approximate
$d(P_{V/d^2}-V/d^2)$ by a sum of $N$ centered $d$-Bernoulli variables,
with success probability $V/(Nd^2)$, for $N$ sufficiently large.
Their total variance is at most $V$ and their upper bound is at most $A$.
The quoted comparison bounds each quadratic hinge by the corresponding
centered $A$-Poisson hinge. Their fourth moments are at most
$d^2V+3V^2$, giving uniform integrability of each quadratic hinge, so the
inequality passes to the Poisson limit. It is preserved by adding an
independent variable: condition first; a bound by a constant times
$V+(Z-y)^2$ justifies integration for finite second moments.
For $w_1<w_2$, split the two comparators into common large-jump variance
$T(1-w_2)$, common small-jump variance $Tw_1$, and the transferred variance
$T(w_2-w_1)$. Replacing the latter small jumps by large jumps gives the
asserted direction. Zero variances are direct.
\end{proof}

The exact own-count variances are
\begin{equation}\label{eq:transport}
 T_0=g_0\Lambda_0,\quad T_1=g_1\Lambda_1,\qquad
 g_0=1-\delta\eta/q,\quad g_1=1+\delta\eta/p.
\end{equation}
For an incoming count, $\delta+i=\xi+r+i>0$. Thus its negative group
has variance at most $\Lambda_0$, while its positive group has variance
at most $\Lambda_1[1+(5/2)(\Xi+r+i)\eta]$.
The inflated comparator at the first incoming count therefore has total
variance at most
\begin{equation}\label{eq:firstT}
 T^+=L\left[1+\frac{5}{84}(\Xi+r_++1)w_+\right],
\end{equation}
and its fraction assigned to small jumps is at least the physical $w$.
For any rational $0\le y<1+r_-$, Lemma~\ref{lem:monotone} and the
pointwise hinge inequality give the whole-cell bound
\begin{equation}\label{eq:first}
 \Pp(W_{k+1}\le-r-1)\le
 \frac{G(T^+,w_-;y)}{(1+r_--y)^2}=:Q.
\end{equation}
One may replace $Q$ by $\min(Q,1)$.

The centered hinge is \emph{not} assumed convex in $w$. Neither do we
replace the own-count allocation by the physical allocation. In general
$T_0/(T_0+T_1)\ne1-w$; the exact factors in \eqref{eq:transport} must
remain in the outgoing moment bound.

\subsection*{An exact finite expression, with no missing Poisson tail}
Put $h=3/125$, so $A=25h$ and $d=3h$. If
$\mu=T(1-w)/A^2$, $\nu=Tw/d^2$, and
$K=y/h+25\mu+3\nu\ge0$, then $G=h^2N$, where
\begin{align}\label{eq:complement}
 N={}&625\mu+9\nu+(25\mu+3\nu-K)^2\\
 &-e^{-\mu-\nu}
 \sum_{\substack{j,k\ge0\\25j+3k\le K}}
       (K-25j-3k)^2\frac{\mu^j\nu^k}{j!k!}.\notag
\end{align}
This subtracts the lower complement from the full second moment. The sum
is finite even for rational $K$, while the identity includes the complete
upper tail. Its coefficients are nonnegative, so a lower exponential bound
produces an upper bound for $N$.
For $x\ge0$, reduce to $z=x/2^a\le1$, sum the first 25 Taylor terms,
and bound the remainder after order 24 by
\[
 \frac{z^{25}/25!}{1-z/26}.
\]
Repeated squaring gives rational endpoints for $e^x$; reciprocals give
endpoints for $e^{-x}$. All arithmetic in acceptance uses integers and
\texttt{Fraction}; no floating-point evaluation enters these bounds.

\subsection*{All later incoming counts}
Every legal conditional count has total variance at most
$n\max(v_0,v_1)\le(5/2)\lambda$. For $i=2,3$ use
\[
 V_i=\min\left(\frac52L,
       L\left[1+\frac5{84}(\Xi+r_++i)\right]\right),\qquad\mu_i=V_i/A^2.
\]
For an integer $j_i\ge2$ with $(i+r_-)/A+\mu_i-j_i>0$, the
quadratic comparison and Poisson factorial moments give
\begin{equation}\label{eq:later}
 \Pp(W_{k+i}\le-r-i)\le
 \frac{\mu_i^{j_i+1}}
 {(j_i+1)!((i+r_-)/A+\mu_i-j_i)^2}=:Q_i.
\end{equation}
Indeed $(P-j)_+^2\le\binom P{j+1}$ for integers $P\ge0,j\ge2$.
For $j=2$ this is $(t-1)(t-2)\ge0$ for integer $t=P-j\ge1$;
the binomial coefficient increases with $j$. Its Poisson expectation is
$\mu^{j+1}/(j+1)!$.
For $i\ge4$, use the common variance $(5/2)L$, put $\mu_0=(5/2)L/A^2$,
and choose hinge $j=i$. The denominator increases with $i$, and the
numerator ratio is at most $\mu_0/6<1$. Hence the entire remainder is
\begin{equation}\label{eq:infinite}
 J=\frac{\mu_0^5}
 {5!\{4(1/A-1)+r_-/A+\mu_0\}^2(1-\mu_0/6)}.
\end{equation}
We take $\mathcal T=Q+Q_2+Q_3+J$. Adding terms for nonexistent counts
can only increase this upper bound. All unsquared hinge denominators are
checked positive before division.

\section{The cubic penalty and the outgoing count use the same moments}
The signs of $Y-p$ on the two conditional supports are fixed, since
$a_0<p$ and $a_1<q$. Put $M_0=nq\E U_0^3$, $M_1=np\E U_1^3$.
Expansion gives
\[
 n\{\E|Y-p|^3-v(1-2v)\}
       =3p\Lambda_0+3q\Lambda_1+M_1-M_0.
\]
The elementary inequality $(1+x)^{3/2}\le1+3x/2+3x^2/8$ implies
\begin{equation}\label{eq:cubic}
 n(v+s)^{3/2}(\rho-\rho_B)
 \ge3p\Lambda_0+3q\Lambda_1+M_1-M_0
               -\frac{39}{50}\lambda-\frac{13}{16}\lambda^2\eta.
\end{equation}
Here $(3/2)(1-2v)\le39/50$ and $3(1-2v)/(8v)\le13/16$.
The right side is positive: it is at least
$\lambda[6/5-d-39/50-(13/16)\lambda\eta]>0$.
Therefore, for $K_0=79/25$, we may use
$C/(v+s)^{3/2}>K_0$ in \eqref{eq:cubic} \emph{before} separating its
signed terms. Also the uniform weaker consequence is
\begin{equation}\label{eq:unitgain}
 C(\rho-\rho_B)>
 \frac{(2/5)\{87/250-(13/16)\lambda\eta\}}{(251/500)^3}
                  \lambda\eta>\lambda\eta.
\end{equation}

For $\alpha>0$ define the quartic minorant
\[
 h_\alpha(x)=\frac{-x^4+3\alpha^2x^2+2\alpha^3x}{5\alpha^4}
                         \le\mathbf1_{(0,\infty)}(x).
\]
Its numerator is $x(x+\alpha)^2(2\alpha-x)$; its positive maximum is
$(9+6\sqrt3)\alpha^4/4<5\alpha^4$ because $108<121$.
For a centered independent sum, expansion gives
\[
 \E W_k^4=\sum\E Z_j^4+3t^2-3\sum(\E Z_j^2)^2,
 \quad\E W_k^3=\sum\E Z_j^3,\quad t=T_0+T_1.
\]
Apply the minorant to $W_k+r$. To combine it with \eqref{eq:cubic},
price each conditional squared variable on its support. For group zero
the price is
\[
 \frac m{5\alpha^4}g_0(x^2+4rx)+K_0\eta x,
 \qquad -A\le x\le a_0\le d;
\]
for group one it is
\[
 \frac m{5\alpha^4}g_1(x^2+4rx)-K_0\eta x,
 \qquad -d\le-a_1\le x\le A.
\]
Both are convex in $x$, because $m,g_0,g_1>0$. Thus their maxima are at
the respective relaxed support endpoints $x_0\in\{-A,d\}$ and
$x_1\in\{-d,A\}$. Dropping the favorable term
$3\sum(\E Z_j^2)^2$ and replacing $3p\ge6/5$, $3q\ge141/100$
gives the following lower bound for \eqref{eq:margin}, for the worst choice
of $x_0,x_1$:
\begin{align}\label{eq:polynomial}
\mathcal L={}&\frac m{5\alpha^4}\big[
 (3\alpha^2-6r^2-x_0^2-4rx_0)T_0
 +(3\alpha^2-6r^2-x_1^2-4rx_1)T_1\notag\\
 &\hspace{27mm}-3(T_0+T_1)^2-r^4+3\alpha^2r^2+2\alpha^3r\big]\notag\\
 &-mr(1+E\eta)-mR\mathcal T-\frac74(\Xi+r)\lambda\eta\\
 &+K_0\eta\lambda\left[
 \frac65(1-w)+\frac{141}{100}w-x_0(1-w)+x_1w
          -\frac{39}{50}-\frac{13}{16}\lambda\eta\right].\notag
\end{align}
This is a same-law moment inequality. The cubic penalty is not allocated
independently of the outgoing fourth-moment estimate.

\section{Why the finite check covers every law and every sample size}
For fixed $r$, the actual multipliers satisfy
\[
 1-\frac52(\Xi+r)\eta\le g_j\le
 1+\frac52(\Xi+r)\eta.
\]
The expression \eqref{eq:polynomial} is separately concave in $g_0,g_1$,
so it suffices to substitute
\[
 g_j=1+\varepsilon_j\frac52(\Xi+r)\eta,
                    \qquad\varepsilon_j\in\{-1,1\}.
\]
This is a valid outer rectangle, even though its four corners need not
be realized by one marginal. After substitution it is separately concave
in each of $\lambda,w,\eta$: all nonaffine terms in any one of those
variables are negative squares. It is affine in $m$. Hence it suffices to
check $\lambda\in\{l,L\}$, $w\in\{w_-,w_+\}$,
$\eta\in\{0,1/42\}$ and $m\in\{m_-,m_+\}$, together with the four
support-endpoint choices. Each choice leaves a polynomial of degree at
most four in $r$. The artificial face $\eta=0$ is an endpoint of a
concavity argument, not a limit theorem for a fixed law.

For a polynomial $\sum_{k=0}^4 c_kr^k$, its degree-four Bernstein
coefficients on $[a,b]$ are
\[
 \beta_i=\sum_{k=0}^4 c_k
  \sum_{j=0}^{\min(i,k)}
   \frac{\binom kj\binom ij}{\binom4j}a^{k-j}(b-a)^j,
                       \qquad0\le i\le4.
\]
This follows by first substituting $r=a+(b-a)t$ and then using
$t^j=\sum_{i=j}^4\binom ij\binom4j^{-1}\binom4i t^i(1-t)^{4-i}$.
Positive coefficients prove positivity on the closed interval.
There are 256 endpoint combinations per quartic cell and five coefficients
per combination.

Some cells use the simpler Cantelli bound. Set
\[
 T=L[1+(5/2)D/42],\quad h_f=1+E/42,\quad
 B=\min_{r\in\{r_-,r_+\}}\left\{\frac{r^2}{T+r^2}-h_fr\right\}
                         -R\mathcal T.
\]
Require $r_->0$ and $2r_-T/(T+r_-^2)^2>h_f$.
The derivative of $r^2/(T+r^2)-h_fr$ first increases and then decreases,
because its derivative has the sign of $T-3r^2$. Under the displayed
left-endpoint condition any interior critical point is a maximum, so the
minimum is at an endpoint. The variance upper bound makes Cantelli
uniform. Using \eqref{eq:unitgain}, a valid margin is
\begin{equation}\label{eq:cantelli}
 \begin{cases}m_-B,&B\ge0,\\m_+B,&B<0\end{cases}
                        -\{(7/4)D-1\}L/42.
\end{equation}
The coefficient $(7/4)D-1$ is positive on the domain.

The certificate \texttt{compound-certificate.json} contains 291 rational
closed boxes, 200 using \eqref{eq:polynomial} and 91 using
\eqref{eq:cantelli}. Its domain is exactly
\[
 [3/100,3/40]\times[0,3/20]\times[0,1].
\]
The initial division has 3 noise intervals, 12 phase intervals and
8 allocation intervals. Three cells were split once, giving 291 final
cells. The acceptance checker independently forms the common endpoint
refinement (480 elementary boxes), checks each interior is covered exactly
once and all exterior faces occur, and uses closed inequalities on every
cell. Missing, duplicated, reordered and malformed cells reject.
Every cell passes with a rational lower margin of at least
\[
                    \frac{256695813}{10^{12}}>0.
\]
This proves Theorem~\ref{thm:comparison} from the stated analytic inputs.
The finite arithmetic is executable evidence for the displayed inequalities;
it does not replace the count identity, comparison theorem or concavity
arguments above.

\section{Scope, checks and preserved failed shortcuts}
The acceptance program is \texttt{compound-accept.py}; it checks 23
normalization gates and 8 malformed-input or source-change controls, reads
the evaluator bytes, verifies their SHA256, and executes those exact bytes
without importing cached bytecode. The finite exponential calculation is
\texttt{compound-evaluator.py}. Its separate controls include zero
intensities, exact hinge collisions, finite positive-series lower bounds,
and malformed floating/Boolean inputs after warming the exact cache.
Floating scripts \texttt{compound-cover-discovery.py} and
\texttt{compound-plan.py} selected the parameters but are not acceptance
inputs. The 294 tested discovery boxes and the first conservative failed
pilot are preserved separately from the exact run.

Two tempting shortcuts are false. First, the physical variance fraction
cannot be reused as the own-count variance fraction. A centered actual law
with $n=42,p=1/2,k=22,r=.08$, $U_0\in\{-.6,.001\}$ and $U_1=-U_0$
has $v_0=v_1=.0006$, $w=1/2$, but own variances $20v_0$ and $22v_1$.
For $\kappa=-4r$ it violates the corresponding shortcut fourth-moment
bound by $423903/2500000000$; its raw $|H|=.92/\sqrt{.2506}<1.84$.
Second, the centered compound hinge need not be convex in allocation:
at $T=1/40,w=1/4,y=1/5$ its second derivative is
\[
 -e^{-815/648}
 \frac{787398879114798488065}{147382613147155571933184}<0.
\]
The independent comparator review derives both controls exactly and supplies
an independent probability-generating-function recurrence for
\eqref{eq:complement}. Neither false shortcut is used here or in the
previous accepted small-noise theorem.

The new comparison can replace the negative-inner part of a selected
right-panel low-noise branch, provided its strip, threshold and noise
hypotheses have been acquired separately and the remaining phases are
covered. It makes half-amplitude confinement and subsequent count-curvature
arguments unnecessary \emph{within this domain}. It does not justify
deleting their uses in other probability panels, larger noise, or other
acquisition and exclusion arguments. No manuscript or inherited certificate
has been changed by this research component.
\end{document}
